\documentclass[aspectratio=169]{beamer}
\usetheme[numbering=fraction, progressbar=frametitle, block=fill]{metropolis}
\usepackage[T2A]{fontenc}
\usepackage[utf8]{inputenc}
\usepackage[russian]{babel}
\usepackage{graphicx}
\usepackage{booktabs}
\usepackage{tikz}
\usepackage{xcolor}

\definecolor{accent}{RGB}{52, 136, 255}
\definecolor{green}{RGB}{102, 221, 136}
\definecolor{orange}{RGB}{255, 204, 68}
\definecolor{red}{RGB}{255, 102, 102}
\definecolor{gray}{RGB}{136, 136, 136}

\setbeamercolor{frametitle}{bg=accent!20, fg=accent}
\setbeamercolor{block title}{bg=accent!20, fg=accent}

\title{Antigravity DEX}
\subtitle{Децентрализованная биржа деривативов}
\date{2026}

\begin{document}

\begin{frame}
\maketitle
\centering{\small PBFT \quad Arbitrum / TON \quad Telegram Mini App}
\end{frame}

\begin{frame}{Проблема: CEX vs DEX}

\begin{columns}[t]
\begin{column}{0.48\textwidth}
\textbf{CEX (Binance, Bybit)}
\begin{itemize}
    \item Деньги на счету компании
    \item Регулятор блокирует счета --- биржа закрыта
    \item Один вход (API Gateway) --- SPOF
    \item Требуют KYC, лицензий
\end{itemize}
\end{column}
\begin{column}{0.48\textwidth}
\textbf{DEX (наша модель)}
\begin{itemize}
    \item Деньги в контрактах/\newline мультисиг-кошельках
    \item Регулятору \textbf{некого блокировать}
    \item Нет единого входа --- 3 независимых сервера
    \item Без KYC --- кошелёк = идентификатор
\end{itemize}
\end{column}
\end{columns}

\vspace{0.5em}
\begin{block}{HyperLiquid --- единственный серьёзный конкурент}
24 валидатора, децентрализованный matching, но:
\textcolor{red}{\textbf{единый CloudFront Gateway}} (2 даунтайма в 2025),
\textcolor{red}{\textbf{\$360k/мес}} инфраструктуры, \textcolor{red}{\textbf{закрытый код}}.
\end{block}
\end{frame}

\begin{frame}{Решение}

\begin{columns}[t]
\begin{column}{0.48\textwidth}
\textbf{Децентрализация без компромиссов}
\begin{itemize}
    \item 3 валидатора, PBFT (2/3 голосов)
    \item Каждый валидатор = свой REST API
    \item Клиент сам выбирает, к кому подключаться
    \item При падении одного --- failover к другому
    \item Никакого API Gateway --- никакого SPOF
\end{itemize}
\end{column}
\begin{column}{0.48\textwidth}
\textbf{Ключевые преимущества}
\begin{itemize}
    \item[\textcolor{green}{\checkmark}] \textbf{Безопасность:} 2/3 подписей для вывода средств
    \item[\textcolor{green}{\checkmark}] \textbf{Open-source:} код публичный, любой может запустить ноду
    \item[\textcolor{green}{\checkmark}] \textbf{Экономия:} \$600/мес vs \$360k/мес (HyperLiquid)
    \item[\textcolor{green}{\checkmark}] \textbf{Регуляторная устойчивость:} \dq мы --- протокол, не биржа\dq
\end{itemize}
\end{column}
\end{columns}
\end{frame}

\begin{frame}{Архитектура (упрощённо)}

\centering
\begin{tabular}{cl}
\textbf{Внешние сети} & \hspace*{1em}\makebox[8em][c]{\textbf{Наша система}} \\
\hline
\makebox[8em][l]{Arbitrum (USDC)} & \hspace*{1em}\makebox[8em][l]{3 валидатора (PBFT)} \\
\makebox[8em][l]{\quad$\cdot$ Rollup batches} & \hspace*{1em}\makebox[8em][l]{\quad$\cdot$ OrderBook \& Margin} \\
\makebox[8em][l]{\quad$\cdot$ Deposits/withdrawals} & \hspace*{1em}\makebox[8em][l]{\quad$\cdot$ Liquidation} \\
\makebox[8em][l]{TON (USDT)} & \hspace*{1em}\makebox[8em][l]{\quad$\cdot$ FlatFileStore} \\
\makebox[8em][l]{\quad$\cdot$ Telegram payments} & \hspace*{1em}\makebox[8em][l]{\quad$\cdot$ Telegram UI} \\
\makebox[8em][l]{\quad$\cdot$ Instant deposits} \\
\end{tabular}

\vspace{0.5em}
\flushleft\small
\textbf{Депозит:} внешняя сеть $\to$ PBFT $\to$ кредит на L3\\
\textbf{Вывод:} PBFT $\to$ внешняя сеть $\to$ пользователь
\end{frame}

\begin{frame}{Сравнение с HyperLiquid}

\centering
\begin{tabular}{lcc}
\toprule
\textbf{Параметр} & \textbf{HyperLiquid} & \textbf{Antigravity DEX} \\
\midrule
Валидаторы & 24 & \textbf{3} \\
Консенсус & HyperBFT & \textbf{PBFT} \\
API Gateway & \textcolor{red}{CloudFront} & \textbf{Нет} \\
TPS & 200k/s & \textbf{$\sim$5M (один валидатор)} \\
Settlement & Свой L1 & \textbf{Arbitrum / TON} \\
Язык & Rust & \textbf{Java 25} \\
Smart contracts & HyperEVM & \textbf{Нет (план)} \\
Native token & HYPE (\$10B) & \textbf{Опционально} \\
Инфраструктура & \textcolor{red}{\textbf{\$360k/мес}} & \textbf{\$600/мес} \\
Код & \textcolor{red}{Закрытый} & \textbf{Open-source} \\
\midrule
\textbf{Регуляторный риск} & \textcolor{orange}{Средний} & \textbf{Низкий} \\
\textbf{Единая точка отказа} & \textcolor{red}{Есть} & \textbf{Нет} \\
\bottomrule
\end{tabular}
\end{frame}

\begin{frame}{Производительность}

\begin{columns}[t]
\begin{column}{0.48\textwidth}
\textbf{Что измерено}
\vspace{0.3em}
\begin{itemize}
    \item[\textcolor{green}{\checkmark}] \textbf{Execution:} $\sim$5M tx/s
    \item In-memory Disruptor, без I/O
    \item ThroughputBenchmarkTest (JVM warmup)
\end{itemize}
\end{column}
\begin{column}{0.48\textwidth}
\textbf{Что предстоит измерить}
\vspace{0.3em}
\begin{itemize}
    \item[\textcolor{orange}{?}] \textbf{Consensus (3 ноды):} TBD
    \item[\textcolor{orange}{?}] \textbf{End-to-end:} TBD
    \item[\textcolor{orange}{?}] \textbf{Latency:} оценка 50--200ms
\end{itemize}
\end{column}
\end{columns}

\vspace{0.8em}
\textbf{Вывод:} узкое место --- не execution, а сетевая задержка PBFT.\\
Consensus throughput: TBD (зависит от объема блока и network latency).
\end{frame}

\begin{frame}{Мосты: Arbitrum + TON}

\begin{columns}[t]
\begin{column}{0.48\textwidth}
\textbf{Arbitrum Bridge}
\begin{itemize}
    \item \textbf{Актив:} USDC
    \item \textbf{Депозит:} 10--15 мин
    \item \textbf{Вывод:} 7 дней / fast 0.3\%
    \item \textbf{Безопасность:} Ethereum finality
    \item \textbf{Challenge window:} 7 дней
\end{itemize}
\end{column}
\begin{column}{0.48\textwidth}
\textbf{TON Bridge}
\begin{itemize}
    \item \textbf{Актив:} USDT
    \item \textbf{Депозит:} $\sim$3 сек
    \item \textbf{Вывод:} $\sim$1 сек (мгновенно)
    \item \textbf{Интеграция:} Telegram Mini App
    \item \textbf{Auth:} WebApp SDK initData
\end{itemize}
\end{column}
\end{columns}

\begin{columns}[t]
\begin{column}{0.48\textwidth}
\textcolor{blue}{\textbf{Arbitrum = L2 Settlement}}\\[0.2em]
\scriptsize
\begin{itemize}
    \item Rollup batch posting (state root + trades)
    \item D\'eпозиты/выводы USDC
    \item 7d challenge window / fast 0.3\%
\end{itemize}
\end{column}
\begin{column}{0.48\textwidth}
\textcolor{orange}{\textbf{TON = Payment Rail}}\\[0.2em]
\scriptsize
\begin{itemize}
    \item \textcolor{gray}{Только д\'eпозиты/выводы, нет батчей}
    \item Telegram Mini App
    \item Instant (\~3s deposit / \~1s withdraw)
\end{itemize}
\end{column}
\end{columns}

\vspace{0.5em}
\begin{block}{\textcolor{orange}{Статус}}
Оба моста --- in-memory mock. Реальная интеграция с Arbitrum Nitro и TON RPC --- первый приоритет для продакшна.
\end{block}
\end{frame}

\begin{frame}{Telegram Mini App}

\textbf{Почему Telegram?}

\begin{itemize}
    \item 900M+ активных пользователей
    \item Встроенный кошелёк (Telegram Wallet --- TON USDT)
    \item WebApp SDK --- авторизация без паролей
    \item Мгновенные депозиты/выводы через TON
\end{itemize}

\textbf{Что сделано:}
\begin{itemize}
    \item UI: telegram.html + telegram-app.js (vanilla JS)
    \item Auth: Telegram WebApp initData (HMAC)
    \item Депозит USDT $\to$ торговля $\to$ вывод USDT
    \item Пользователи tg-* пропускают RSA-верификацию
\end{itemize}

\textbf{План:}
\begin{itemize}
    \item TON Connect / Tonkeeper интеграция
    \item Push-уведомления о сделках
    \item In-app графики
\end{itemize}
\end{frame}

\begin{frame}{Бизнес-модель}

\begin{columns}[t]
\begin{column}{0.48\textwidth}
\textbf{Источники дохода}
\begin{itemize}
    \item \textbf{Protocol fee:} 0.02--0.05\% с каждой сделки
    \item Fast withdrawal fee: 0.3\% (LP + foundation)
    \item Funding rate: спред забирает протокол
    \item Native token (будущее): стейкинг, DAO, fee discount
\end{itemize}
\end{column}
\begin{column}{0.48\textwidth}
\textbf{Расходы}
\begin{itemize}
    \item 3 $\times$ bare metal: $\sim$\$150/мес
    \item Arbitrum batches: $\sim$\$17--50/мес
    \item TON RPC: $\sim$\$0 (public)
    \item Домены, SSL: $\sim$\$0
    \item \textbf{Итого: $\sim$ \$200/мес}
\end{itemize}
\end{column}
\end{columns}

\vspace{0.5em}
\begin{block}{Считаем прибыль}
\begin{tabular}{lr}
\textbf{Дневной объём} & \textbf{Доход (0.03\% fee)} \\
\hline
\$100k & \$30/день или \textasciitilde\$900/мес \\
\textbf{\$1M} & \textbf{\$300/день или \textasciitilde\$9k/мес} \\
\$10M & \$3k/день или \textasciitilde\$90k/мес \\
\hline
\multicolumn{2}{l}{\small\textcolor{gray}{Для сравнения: HyperLiquid $\sim$\$2--5B/день}} \\
\end{tabular}
\end{block}

\flushleft\small
При объёме \$1M/день --- чистая прибыль \textbf{\$8.8k/мес} (после расходов).
\end{frame}

\begin{frame}{Регуляторная стратегия}

\begin{enumerate}
    \item \textbf{Cayman foundation} владеет IP и торговой маркой, \textbf{но не управляет биржей}
    \item \textbf{Open-source} --- код публичный. Нельзя запретить протокол --- можно запретить конкретную компанию, но у нас её нет
    \item \textbf{Без KYC} --- кошелёк = идентификатор. Мы физически не собираем персональные данные
    \item \textbf{Геоблокировка US} (Cloudflare IP) --- на уровне UI, не на уровне протокола. Валидаторы в других юрисдикциях продолжают работать
    \item \textbf{Прецедент HyperLiquid}: Cayman + open-source + без KYC, \$2--5B/день, SEC не может закрыть --- потому что закрывать некого
\end{enumerate}

\vspace{0.5em}
\begin{block}{Ключевое сообщение}
Мы не биржа в юридическом смысле. Мы --- протокол (набор правил и кода),
через который пользователи торгуют друг с другом. Валидаторы --- не сотрудники,
а независимые операторы. Как майнеры в Bitcoin.
\end{block}

\centering{\small \textcolor{orange}{Не является юридической консультацией. Требуется advice от crypto-юриста.}}
\end{frame}

\begin{frame}{Roadmap}

\begin{columns}[t]
\begin{column}{0.48\textwidth}
\textcolor{green}{\textbf{Сделано}}
\begin{itemize}
    \item PBFT consensus (3 валидатора, 2/3)
    \item LMAX Disruptor pipeline
    \item OrderBook + MarginManager
    \item Liquidation + Funding
    \item ArbitrumBridge + TonBridge (mock)
    \item FlatFileStore (WAL + snapshot)
    \item Telegram Mini App
    \item 36 тестов, все проходят
\end{itemize}
\end{column}
\begin{column}{0.48\textwidth}
\textcolor{orange}{\textbf{Следующие шаги}}
\begin{enumerate}
    \item \textbf{Currency model:} oracle-weighted equity
    \item \textbf{Real bridges:} Arbitrum Nitro + TON RPC
    \item \textbf{TON Vault:} multisig / smart contract
    \item \textbf{Validator economics:} stake, slashing, fees
    \item Network benchmark (3 региона)
    \item Security audit
\end{enumerate}
\vspace{0.5em}
\textbf{Детали:} github.com/rgobov/dex-on-java/UNRESOLVED.md
\end{column}
\end{columns}
\end{frame}

\begin{frame}{Что нужно сейчас}

\begin{columns}[t]
\begin{column}{0.33\textwidth}
\centering
\textbf{1. Security audit}\\[0.3em]
\small
PBFT consensus, bridges, margin logic, cryptography.\\
Базово --- до реальных денег.
\end{column}
\begin{column}{0.33\textwidth}
\centering
\textbf{2. Real bridge integration}\\[0.3em]
\small
Arbitrum Nitro (RPC + retryable tickets) + TON RPC.\\
Замена mock на реальные вызовы.
\end{column}
\begin{column}{0.33\textwidth}
\centering
\textbf{3. Validator economics}\\[0.3em]
\small
Stake model, slashing conditions, fee distribution.\\
Как привлечь независимых валидаторов.
\end{column}
\end{columns}

\vspace{1em}
\centering
\large\textbf{Готовы к диалогу с инвесторами и техническими партнёрами.}

\vspace{0.5em}
\small\textcolor{gray}{github.com/rgobov/dex-on-java}
\end{frame}

\end{document}
